\documentclass[11pt,a4paper]{article}
\usepackage[T1]{fontenc}
\usepackage[utf8]{inputenc}
\usepackage{lmodern}
\usepackage[margin=2.5cm]{geometry}
\usepackage{booktabs}
\usepackage{tabularx}
\usepackage{enumitem}
\usepackage{listings}
\usepackage{xcolor}
\usepackage[hidelinks]{hyperref}

\lstset{
  basicstyle=\ttfamily\small,
  columns=fullflexible,
  keepspaces=true,
  frame=single,
  rulecolor=\color{gray!50},
  backgroundcolor=\color{gray!8},
  breaklines=true,
  upquote=true,
  showstringspaces=false,
  xleftmargin=0.5em,
  xrightmargin=0.5em
}
\setlength{\emergencystretch}{3em}
\newcommand{\kbd}[1]{\texttt{#1}}
\newcolumntype{L}{>{\raggedright\arraybackslash}X}

\title{LeRobot (Free Pascal) Installation Manual\\[0.3em]\large Installing LeRobot on another Mac}
\author{}
\date{9 November 2026}

\begin{document}
\maketitle

This manual explains how to install LeRobot for Free Pascal on a Mac where
it has never been installed. With it, the Turbo Pascal programs of the 1991
book \emph{Premières leçons de programmation en Turbo Pascal} are compiled
by Free Pascal and run in the robot window. Once it is done, use the
\emph{LeRobot (Free Pascal) User Manual} (\kbd{LeRobot\_Manual.pdf}) to draw
and to run robot programs.

The robot window is the one of the Python version (an ordinary macOS window
drawn with tkinter): the Pascal programs send their commands to it. So you
need three things: the Free Pascal compiler, a Python with Tk 8.6, and the
\kbd{LeRobot} folder with its \kbd{LeRobot\_FreePascal\_FR} \emph{and}
\kbd{LeRobot\_Python\_FR} folders. No XQuartz is needed. The installation takes
about twenty minutes. You need:
\begin{itemize}
\item a Mac with a recent macOS (Apple Silicon or Intel);
\item an administrator account (the installers ask for the Mac password);
\item an Internet connection;
\item a copy of the \kbd{LeRobot} folder (USB key, AirDrop, cloud drive,
  \dots).
\end{itemize}

Commands are typed in the \textbf{Terminal} application
(\emph{Applications} $\rightarrow$ \emph{Utilities} $\rightarrow$
\emph{Terminal}). Copy each command in one piece, paste it into the
Terminal, and press Return.

\tableofcontents

%======================================================================
\section{The Free Pascal compiler}

\subsection{Command Line Tools}
Free Pascal uses the linker of Apple's Command Line Tools. If Xcode or the
Command Line Tools are not installed yet:
\begin{lstlisting}[language=bash]
xcode-select --install
\end{lstlisting}
and accept the installation in the window that opens (a few minutes).

\subsection{Free Pascal with Homebrew}
Homebrew (\url{https://brew.sh}) is the simplest way to install Free
Pascal. If it is not installed, install it with the command shown on its
home page, then:
\begin{lstlisting}[language=bash]
brew install fpc
\end{lstlisting}
Check, in a new Terminal window:
\begin{lstlisting}[language=bash]
fpc -iV
\end{lstlisting}
The answer is the version of Free Pascal (\kbd{3.2.2} or later).

Free Pascal can also be installed without Homebrew, with the installer of
\url{https://www.freepascal.org} (\emph{Download}, \emph{macOS}).

\paragraph{Messages to ignore.} On recent macOS, Free Pascal 3.2.2 prints
linker warnings such as \kbd{-macosx\_version\_min has been renamed} or
\kbd{-multiply\_defined is obsolete}. They are harmless; \kbd{robot.sh}
hides them.

%======================================================================
\section{Python with Tk 8.6}

The robot window needs Python 3.9 or later \emph{linked to Tk 8.6 or later}.
Tk is the library that draws the window.

\textbf{The \kbd{python3} that comes with macOS is not enough.} It is
linked to Tk 8.5, which no longer works on recent macOS: the robot window
does not open and the program hangs. Install Python from python.org, whose
installer includes a recent Tk:
\begin{enumerate}
\item Open \url{https://www.python.org/downloads/macos/} and download the
  \emph{macOS 64-bit universal2 installer} of the latest Python 3 release.
\item Open the downloaded \kbd{.pkg} file and follow the installer. At the
  end, the Finder window it opens can be closed.
\end{enumerate}
Check, in a new Terminal window:
\begin{lstlisting}[language=bash]
/Library/Frameworks/Python.framework/Versions/Current/bin/python3 -c "import tkinter; print(tkinter.TkVersion)"
\end{lstlisting}
The answer must be \kbd{8.6} or more. You do not need to change which
\kbd{python3} the Terminal uses: when the robot window is started with the
macOS \kbd{python3}, it notices that its Tk is too old and starts itself
again with the python.org Python.

%======================================================================
\section{The LeRobot folder}

Copy the \kbd{LeRobot} folder anywhere, for example onto the
\textbf{Desktop}. The Pascal programs are in its folder
\kbd{LeRobot\_FreePascal\_FR}; the folder \kbd{LeRobot\_Python\_FR} must stay next
to it, in the same \kbd{LeRobot} folder: it draws the robot window. The
\kbd{LeRobotFreePascal.app} icon must stay in \kbd{LeRobot\_FreePascal\_FR},
next to \kbd{robot.sh}. The whole \kbd{LeRobot} folder can be moved or
renamed later: nothing in it depends on where it is.

\paragraph{Commands.} The \kbd{LeRobot} folder can be anywhere (on the
Desktop, in Documents, \dots): everything in it is found relative to it.
The commands of this manual are typed in a Terminal window \emph{opened in
the \kbd{LeRobot} folder}: type \kbd{cd} and a space, drag the
\kbd{LeRobot} folder from the Finder onto the Terminal window, and press
Return (or, in the Finder, right-click the folder and choose
\emph{Services} $\rightarrow$ \emph{New Terminal at Folder}).

If the folder comes from another Mac, make the scripts executable again (a
copy can lose this) and remove the compiled files of the other Mac (they
are made again when needed):
\begin{lstlisting}[language=bash]
cd LeRobot_FreePascal_FR && chmod +x robot.sh verifier.sh outils/*.sh outils/app/*.sh && rm -rf _build
\end{lstlisting}

%======================================================================
\section{First launch}

\begin{enumerate}
\item Run a program of the book:
\begin{lstlisting}[language=bash]
cd LeRobot_FreePascal_FR
./robot.sh livre/03_programmation/Carre.p
\end{lstlisting}
  The first time, \kbd{robot.sh} compiles the robot units (a few seconds),
  then the program. The robot window ($800 \times 700$ points) opens in
  front of the other windows. When the square is drawn, press Esc to close
  the window.
\item Open the interactive drawing program:
\begin{lstlisting}[language=bash]
./robot.sh livre/Didactitiels_et_jeux/Dessiner.p
\end{lstlisting}
  Press Return in the window to finish, then \kbd{n} not to save.
\item \kbd{LeRobotFreePascal.app}, in the \kbd{LeRobot\_FreePascal\_FR}
  folder: double-click its icon to run Dessiner; drop a \kbd{.p} program on
  its icon to compile and run it. The program runs in a Terminal window,
  which closes when the program ends normally and stays open if there is
  an error (of compilation or of execution).

  The first time you open \kbd{LeRobotFreePascal.app}, macOS may refuse to
  open it because it comes from another Mac. Right-click the icon, choose
  \emph{Open}, then confirm. Allow it to control Terminal and to access the
  folder of \kbd{LeRobot} (Desktop, Documents, \dots) when macOS asks.
\item If the folder was copied from another Mac and the application still
  refuses to open, rebuild it on this Mac:
\begin{lstlisting}[language=bash]
cd LeRobot_FreePascal_FR && outils/app/fabriquer.sh
\end{lstlisting}
\end{enumerate}

%======================================================================
\section{Troubleshooting}

\begin{tabularx}{\textwidth}{@{}LLL@{}}
\toprule
\textbf{Symptom} & \textbf{Cause} & \textbf{Fix} \\
\midrule
\emph{command not found: fpc} & Free Pascal is not installed, or the Terminal was opened before & \kbd{brew install fpc} (Section~1); open a new Terminal window \\
\addlinespace
\emph{Error: ld: \dots} or \emph{linker not found} & No Command Line Tools & \kbd{xcode-select --install} (Section~1) \\
\addlinespace
\emph{Can't find unit Robot1} & The robot units are not compiled (folder \kbd{\_build} removed during compilation) & Run \kbd{./robot.sh} again: it compiles them \\
\addlinespace
The program hangs and no window opens & Python linked to Tk 8.5 (the macOS \kbd{python3}) and no other Python found & Install Python from python.org (Section~2) \\
\addlinespace
\emph{can't open file '.../serveur.py'} or \emph{No module named 'LibRobot'} & The folder \kbd{LeRobot\_Python\_FR} is missing or not next to \kbd{LeRobot\_FreePascal\_FR} & Copy the whole \kbd{LeRobot} folder (Section~3) \\
\addlinespace
\emph{permission denied: ./robot.sh} & The scripts lost their executable mark during the copy & \kbd{chmod +x robot.sh verifier.sh} (Section~3) \\
\addlinespace
\emph{Ce programme doit etre lance par robot.sh} printed by a program & The compiled program was started by itself, without the robot window & Use \kbd{./robot.sh Program.p} or \kbd{LeRobotFreePascal.app} \\
\addlinespace
\kbd{LeRobotFreePascal.app} does not open & macOS blocks applications copied from another Mac & Right-click, \emph{Open}; or rebuild it (Section~4) \\
\addlinespace
To use a given Python & --- & \kbd{LEROBOT\_PYTHON=\allowbreak /path/\allowbreak to/\allowbreak python3 \allowbreak ./robot.sh Program.p} \\
\bottomrule
\end{tabularx}

\end{document}
